\documentclass[11pt]{article}
\usepackage[utf8]{inputenc}
\usepackage[T1]{fontenc}
\usepackage{lmodern}
\usepackage[margin=1in]{geometry}
\usepackage{csquotes}
\usepackage[style=authoryear,backend=biber,sorting=nyt,maxbibnames=99,giveninits=true,uniquename=init]{biblatex}
\usepackage[hidelinks]{hyperref}
\addbibresource{references.bib}
\begin{document}
\section*{Bioaugmentation to Enhance Woodchip Bioreactor Denitrification}

\subsection{Overview}
Woodchip denitrification bioreactors remove nitrate from agricultural subsurface drainage through microbially mediated denitrification, yet their performance remains constrained by low water temperatures, variable hydraulic retention times (HRTs), and the biogeochemical evolution of the woodchip medium as it ages. A central open question is whether the microbial engine of these systems can be deliberately enhanced---for example by \emph{bioaugmentation}, the inoculation of cold-adapted or otherwise advantageous denitrifying microorganisms---to raise nitrate removal beyond what design and management refinements alone can achieve. This frontier is unsettled. Some evidence suggests that inoculating specific denitrifiers might increase removal rates under low-temperature conditions \parencite{Jang2019}, and dedicated confirmatory work is planned to test this in the field \parencite{Anderson2020}. Other work counsels caution, arguing that optimization is more likely to come from refinements in design and management parameters that shape environmental conditions than from direct manipulation of the microbial community \parencite{Hathaway2015}. Resolving this tension requires progress along several interlocking research lines: (i) confirming the field efficacy, survival, and persistence of inoculated strains and clarifying their reduction pathways to dinitrogen (N\textsubscript{2}); (ii) building a mechanistic understanding of the resident bioreactor microbiome, its substrate metabolism, and its production of harmful by-products; (iii) resolving the hydraulics, HRT dependence, and carbon dynamics that govern removal; (iv) characterizing long-term aging and longevity; and (v) integrating bioreactors with complementary practices, novel water sources, and enhanced treatment configurations. The findings synthesized below are drawn overwhelmingly from site-scale field and laboratory studies concentrated in the U.S. Midwest (Iowa, Minnesota, South Dakota) with additional evidence from the Mid-Atlantic, Ohio, Virginia, North Carolina, Lithuania, Japan, and global syntheses, and they collectively define an agenda framed around what remains to be tested rather than what is already established.

\subsection{Bioaugmentation and Direct Microbial Manipulation}
The defining frontier of this theme is whether denitrifying microorganisms can be inoculated into woodchip bioreactors to enhance nitrate removal, especially where cold temperatures limit biological activity. \textcite{Jang2019} propose that inoculating denitrifiers (bioaugmentation) might increase the nitrate removal rate of woodchip bioreactors at low-temperature conditions, but they note that genome completion and further data mining are still necessary to identify the nitric oxide (NO) reductase of \emph{Cellulomonas} sp.\ strain WB94---an unresolved step in confirming a complete denitrification pathway for a candidate inoculant.

\textcite{Anderson2020} frame the key unanswered questions explicitly: further research is necessary to confirm the efficacy of inoculating these denitrifiers for enhancing nitrate removal in woodchip bioreactors under field conditions, including whether inoculated strains actually enhance nitrogen removal in the field, how long inoculated strains can survive in bioreactors, and whether carbon addition can enhance the abundance of both the inoculated and the indigenous denitrifiers; the authors state they are planning to address these questions in the near future. Survival and persistence of the inoculum thus emerge as first-order uncertainties, echoed at the global scale by \textcite{Wang2025b}, who call for developing field-scale bioaugmentation methods that minimize inoculum loss and ensure persistence of inoculated cultures, and for expanding knowledge of \emph{nod}-harboring bacteria across different agricultural settings.

A parallel line concerns the terminal steps of denitrification and the microbes that catalyze them. \textcite{Aldossari2021} report that further research is necessary to confirm the reduction of nitrous oxide (N\textsubscript{2}O) to N\textsubscript{2} in fungal strains, including identification of nitrous oxide reductase genes on their genomes, and---critically for the theme---to see whether these strains can reduce nitrate to N\textsubscript{2} under woodchip bioreactor conditions. They add methodological caveats: confirmatory experiments are needed to examine the usefulness of the acetylene inhibition technique for studying fungal N\textsubscript{2}O reduction, and more defined media should be used in future fungal denitrification studies to clarify the nitrate-consuming pathways.

Beyond conventional heterotrophic denitrifiers, \textcite{Wang2025b} identify emerging microbial targets and strategies for future work: exploring the potential of bioaugmenting methanogenic environments with nitrate-dependent anaerobic methane oxidation (n-DAMO) and other oxygenic denitrifiers; providing suitable conditions as a biostimulation strategy to enhance the competitiveness of inoculated n-DAMO bacteria against other nitrogen-cycling processes; and exploring algal-bacterial granular sludge technology for agricultural drainage treatment. These directions extend bioaugmentation from single-strain inoculation toward engineered consortia and coupled metabolic guilds.

Against this enthusiasm, \textcite{Hathaway2015} offer a counter-hypothesis worth testing directly: that optimization of denitrifying bioreactors is likely to result from refinements in design and management parameters influencing environmental conditions rather than through inoculation or other direct manipulations of the microbial community. Whether bioaugmentation can outperform---or must be combined with---such environmental engineering remains an open, testable question and a useful null hypothesis for future field trials.

\subsection{Understanding the Resident Microbiome and By-Product Formation}
Rational manipulation of bioreactor microbiology presupposes a mechanistic understanding of the resident community, which remains incomplete. \textcite{Hartfiel2022} argue that future work should draw on the knowledge of soil and wetland ecology to inform the study of bioreactor microbiomes, and that further study is warranted of the factors influencing cellulose and lignin metabolism genes and the community members that promote substrate degradation. They also highlight processes beyond canonical denitrification: dissimilatory nitrate reduction to ammonium (DNRA) and anammox have been observed in bioreactor systems and warrant study given the limited research into them, and varying substrates to select for differing microbial communities is a promising avenue to promote consistency and potentially reduce by-product formation.

By-product and pollution-swapping risks are a recurring frontier. \textcite{Hartfiel2022} note that the role of the microbial community in harmful by-product formation---including greenhouse gas (GHG) production, sulfate reduction, and subsequent methylmercury production---warrants study, and that although HRT has been implicated as a design factor contributing to by-product production, the mechanism by which HRT impacts the microbial community should be examined. Complementing this, \textcite{Bock2018} suggest that continued study of GHG emissions might enable a more precise comparison of the nitrate-nitrogen removal efficiency gained by increasing HRTs against the additional methane (CH\textsubscript{4}) produced, and that managing bed pH may be important for applications with acidic agricultural drainage to enhance nitrogen removal and mitigate N\textsubscript{2}O emissions. \textcite{Bock2018} further propose that conceptualizing how regional differences impact in-bed controls on nitrogen removal will guide adaptation of bioreactor designs to a wider range of agroecosystems. At the operational level, \textcite{Kouanda2021} recommend that nitrite be evaluated during woodchip bioreactor operation to avoid high nitrite levels in reactor effluents---a specific pollution-swapping pathway needing routine monitoring. \textcite{Wang2025b} generalize the concern, urging that future efforts optimize carbon utilization efficiency to minimize pollution swapping and reduce carbon consumption by non-denitrifying processes, and that research clarify the long-term effects of various conservation practices on the soil and sediment microbial communities responsible for nitrate remediation.

Selective pressures on community assembly also merit dedicated study. \textcite{Schaefer2021} propose characterizing the effect of oxygen load on microbial community composition in contrast to the selective pressure of nitrate load, and argue that further research on the microbial characteristics of denitrifying bioreactors is warranted to explore differences in decomposition and denitrification performance across bioreactors operating at different HRTs.

\subsection{Hydraulics, Hydraulic Retention Time, and Design Guidelines}
HRT is repeatedly identified as the master variable linking hydraulics to removal, yet its treatment in design guidance remains imperfect. \textcite{Moorman2015} state that increased research on the relationship between HRT and nitrate-nitrogen removal would improve bioreactor sizing criteria, while \textcite{Povilaitis2018} call for better optimization and proper evaluation of the effects of HRT to improve the design of denitrifying woodchip bioreactors, and for more investigations at both laboratory and field scales to quantify long-term performance and optimize nitrate-nitrogen removal rates for various flow conditions. In a global synthesis, \textcite{Christianson2025c} recommend several methodological advances: assessing varying HRT by changing flow rate versus sampling along the reactor length; increasing transparency in methods for estimating woodchip porosity; reporting parameters such as dissolved oxygen, pH, and dissolved organic carbon (DOC) more consistently; using machine learning approaches as more data become available to refine HRT--nitrate-removal-efficiency parameter estimates; and continuing work across all scales and bioreactor ages to refine the science and engineering of nitrogen removal.

The internal hydraulics of woodchip media are themselves poorly resolved and warrant flow-physics research. \textcite{Ghane2016} argue that future research should use the non-Darcy flow characteristics of woodchips to adjust the current NRCS design guidelines for denitrification beds and yield more efficient systems. Tracer methodology is a specific bottleneck: \textcite{Ghane2019b} raise concerns about low bromide recovery in denitrification beds and recommend a suite of investigations---testing for organobromine in woodchips following bromide tracer tests, examining the interior pores of woodchips under a microscope to assess bromide entrapment, excavating a woodchip bed after tracer testing to analyze woodchips for physically trapped bromide, ascertaining the cause of low bromide recovery, and either accounting for possible retardation of bromide to woodchips or using non-sorbing conservative tracers such as isotopes to investigate internal hydraulics. \textcite{Ghane2018} likewise recommend conducting tracer tests to determine differences in hydraulic performance of retrofitted beds so that beds can be compared in subsequent experiments. Together these point to an unresolved need for reliable tracer methods before hydraulic performance can be benchmarked.

\subsection{Carbon Dynamics, Substrate Amendment, and Startup}
Because woodchip carbon fuels denitrification, its release kinetics, biodegradability, and pre-treatment are active frontiers. \textcite{Kouanda2021} report that long-term laboratory and field studies are needed to evaluate the longevity of the effect of woodchip composting on nitrate removal, and that more studies are needed to determine the biodegradability of woodchip DOC and the release rates of biodegradable organic carbon to better understand nitrate removal kinetics. Startup carbon losses raise both water-quality and downstream-treatment concerns: \textcite{Abusallout2017} suggest that the initial flush of woodchips may be recycled back to agricultural fields to reduce high organic carbon loadings during startup, that pre-washing woodchips before installation could reduce the initial DOC-releasing potential, and that drinking water treatment plants may need to adjust coagulant types and doses to effectively remove woodchip-derived DOC and reduce disinfection by-product (DBP) formation potential. Substrate selection, as noted above, is also proposed by \textcite{Hartfiel2022} as a lever to shape microbial communities and by-product outcomes.

\subsection{Aging, Longevity, and Long-Term Performance}
The multi-decadal service life of bioreactors is one of the least-studied dimensions and a clear research priority. \textcite{Ghane2018} recommend continuing investigation of the reduction of woodchip particle size over a longer period because particle size is an important factor affecting water flow through denitrification beds; investigating whether the woodchip carbon becomes recalcitrant to the point of nitrate-removal inefficiency before particle-size reduction affects flow; and conducting longer-term studies beyond ten years to better understand how age affects woodchip carbon properties and nitrate removal performance. \textcite{Christianson2020} similarly call for long-term studies documenting biofilm formation and flushing, siltation, wood decay, and additional water chemistry parameters across the life of a bioreactor to better inform threshold values for woodchip recharge and guide assessment of aging bioreactors, together with additional research on the effects on denitrification performance of decreasing carbon-to-nitrogen (C:N) ratio and increasing lignocellulose index (LCI) over time. Long-term field programs to establish efficacy are ongoing, with \textcite{Greenan2009} noting that long-term ($>$8 yr) field studies were underway to assess the long-term efficacy of woodchip denitrification bioreactors. Spatiotemporal monitoring under variable hydrology is also flagged: \textcite{Hartfiel2023b} recommend evaluating spatiotemporal nitrate-nitrogen removal under normal, wet, and drought years and across longer time periods, and combining spatiotemporal nitrate monitoring with woodchip C:N ratio measurements to identify strategies to improve bioreactor design.

\subsection{Management, Novel Water Sources, and Coupled Practices}
A growing frontier concerns active management and the extension of bioreactors beyond conventional gravity-fed drainage. \textcite{Hartfiel2023} advance recommendations for pumped bioreactor systems: using the actual bioreactor system capacity rather than the design capacity for sizing future pumping systems; timing pumping to consider source-water nitrate concentration, temperature, and carbon content; avoiding pumping during periods of high temperatures and/or carbon levels with lower nitrate concentrations; and extending pumped applications to nitrate-laden surface waters, irrigation waters, drainage ditches, pumped cistern systems for stored subsurface drainage, or groundwater. \textcite{Hartfiel2023b} propose additional management experiments: evaluating partial bioreactor recharge strategies near the inlet and their long-term effects; developing similar management recommendations for bioreactors treating alternative nitrate-laden water sources; and managing dual-chamber bioreactors so that only one parallel chamber is used under traditional or low-flow conditions while the second is engaged under higher-flow conditions.

Novel source waters raise their own questions. \textcite{Easton2019} note that additional research is needed to evaluate how denitrification rates in bioreactors respond to the flow and concentration conditions unique to springs under field conditions, and that water policy changes---including financial incentive systems based on payments for quantities of nitrogen removed---may be needed to incentivize landowner adoption of spring bioreactors and stimulate investment in treating legacy nitrogen. At the broadest scale, \textcite{Wang2025b} emphasize that effective management of agricultural subsurface drainage will likely depend on integrating multiple conservation practices to achieve targeted nitrate-reduction goals, and that future studies should clarify the effects of controlled drainage on the denitrifying microbiome and the variables affecting its performance as a nitrogen-mitigation method.

\subsection{Enhanced and Coupled Treatment Configurations}
Beyond passive woodchip beds, several enhanced configurations define emerging research directions. \textcite{Law2018} recommend testing the performance of a well-designed bioelectrochemical reactor (BER) under low-temperature conditions and propose specific design refinements---a horizontal-flow reactor with an upstream anode, baffles to separate oxidizing and reducing zones, and plate-shaped cathodes to improve denitrification---together with lower current intensity and larger cathode surface area as key changes to reduce nitrate-removal cost and improve performance. For co-treatment of phosphorus, \textcite{Penn2024} recommend constructing phosphorus removal structures at phosphorus critical source areas for maximum efficiency, designing bioreactors and phosphorus removal structures to treat as much peak flow as possible, and pursuing further research to confirm or reject the proposed redox-coupled Fe(II)--nitrate reaction for nitrate removal by activated aluminum. These coupled systems broaden the design space in which enhanced denitrification---whether biological or electrochemical---might be pursued.

\subsection{Cross-Cutting Gaps and Methodological Priorities}
Several open questions cut across the theme and constrain progress on the others. The mechanisms of nitrate export from agricultural fields and of denitrification in woodchip bioreactors require further investigation to generate practical guidelines \parencite{Liu2020}, and more events with different durations are needed for robust analysis of the influence of antecedent conditions on hysteresis patterns, as are investigations of the legacy effects of animal-waste application and its interaction with plant growth, temperature dynamics, and drying--wetting periods \parencite{Liu2020}. The fate of organic chemicals is also unresolved: removal of organic chemicals by woodchip bioreactors needs to be investigated under field conditions \parencite{Ilhan2011}. Finally, complementary in-situ interception methods warrant testing, as \textcite{Pampolino2000} suggest that installing a sufficient number of resin capsules in situ (for example, 10 capsules per 135 m\textsuperscript{2}) may allow interception and adsorption of a representative amount of nitrate transported via macropore or even bypass flow. Collectively, these cross-cutting priorities---mechanistic field studies, event- and legacy-scale monitoring, contaminant fate, and improved in-situ measurement---underpin the more targeted frontier of enhancing bioreactor denitrification through bioaugmentation and coupled design.

\printbibliography
\end{document}
